\documentclass[11pt,a4paper]{article}
\usepackage{pmo_activity_handout}

\hypersetup{colorlinks=true,urlcolor=pmoTeal,pdftitle={Dashboard Diagnostic Practice Lab},pdfauthor={Mohamad Aoude}}
\rhead{Practice Exercise}

\begin{document}

\activitytitle{PRACTICE ACTIVITY}{Dashboard Diagnostic Lab}{Calculate, Interpret, Decide}{Duration: 30 minutes \quad | \quad Work: individual with feedback \quad | \quad Output: completed KPI diagnosis}

\begin{activitybox}
\textbf{Practice mission.} Use the week 6 dashboard to calculate project-performance indicators, interpret what each result means, and select a corrective response. Complete each checkpoint before consulting the self-check page.
\end{activitybox}

\section{1. Learning Outcomes}

Students will calculate schedule and cost indicators, interpret their direction and magnitude, connect quantitative and quality evidence, and select corrective action with immediate self-feedback.

\section{2. Week 6 Dataset}

\begin{center}
\begin{tabularx}{0.9\textwidth}{>{\bfseries}Y >{\centering\arraybackslash}p{3cm}}
\toprule
Measure & Value \\
\midrule
Budget at Completion (BAC) & USD 40,000 \\
Planned progress & 50\% \\
Actual completed value & 38\% \\
Actual Cost (AC) & USD 24,000 \\
Planned modules completed & 4 \\
Actual modules completed & 2 \\
Critical bugs open & 12 \\
\bottomrule
\end{tabularx}
\end{center}

\section{3. Formula Reference}

\begin{center}
\begin{tabularx}{\textwidth}{>{\bfseries}p{3cm} p{4.5cm} Y}
\toprule
Indicator & Formula & Interpretation \\
\midrule
Planned Value (PV) & BAC $\times$ planned progress & Budgeted value of work planned \\
Earned Value (EV) & BAC $\times$ actual completed value & Budgeted value of work completed \\
Schedule Variance (SV) & EV $-$ PV & Negative means behind schedule \\
Cost Variance (CV) & EV $-$ AC & Negative means over budget \\
Schedule Performance Index (SPI) & EV $\div$ PV & Below 1 means inefficient schedule performance \\
Cost Performance Index (CPI) & EV $\div$ AC & Below 1 means inefficient cost performance \\
\bottomrule
\end{tabularx}
\end{center}

\section{4. Checkpoint A -- Calculate PV and EV (5 minutes)}

Show the calculation and include units.

\textbf{PV =} \blankline

\vspace{6mm}
\textbf{EV =} \blankline

\vspace{6mm}
In one sentence, explain why PV and EV are expressed as monetary values even though the inputs are percentages:

\vspace{8mm}\blankline

\section{5. Checkpoint B -- Calculate Variances (5 minutes)}

\textbf{SV = EV $-$ PV =} \blankline

\vspace{6mm}
\textbf{CV = EV $-$ AC =} \blankline

\vspace{6mm}
Mark and justify the diagnosis:

\begin{itemize}
  \item Schedule: $\square$ ahead \quad $\square$ on plan \quad $\square$ behind
  \item Cost: $\square$ under budget \quad $\square$ on plan \quad $\square$ over budget
\end{itemize}

\textbf{Evidence-based explanation:}

\vspace{8mm}\blankline

\section{6. Checkpoint C -- Calculate Performance Indices (5 minutes)}

Round answers to two decimal places.

\textbf{SPI = EV $\div$ PV =} \blankline

\vspace{6mm}
\textbf{CPI = EV $\div$ AC =} \blankline

\vspace{6mm}
Complete the interpretations:

For each USD 1.00 of planned work by week 6, the project completed approximately USD \underline{\hspace{2cm}} of value.

For each USD 1.00 spent, the project produced approximately USD \underline{\hspace{2cm}} of earned value.

\section{7. Checkpoint D -- Integrate Non-Financial Evidence (5 minutes)}

Quantitative indicators do not explain everything. Connect each additional signal to a project concern.

\begin{center}
\begin{tabularx}{\textwidth}{>{\bfseries}p{4cm} Y Y}
\toprule
Signal & Concern indicated & Question for the PMO \\
\midrule
2 of 4 modules completed & & \\[7mm]
12 critical bugs open & & \\[7mm]
Database integration risk & & \\
\bottomrule
\end{tabularx}
\end{center}

\section{8. Checkpoint E -- Decide (5 minutes)}

Rank these actions from 1 (strongest) to 4 (weakest):

\begin{itemize}
  \item[$\square$] Add optional features to improve sponsor satisfaction.
  \item[$\square$] Freeze non-critical scope and focus on integration, essential modules, and critical bugs.
  \item[$\square$] Stop reporting until the team recovers lost progress.
  \item[$\square$] Remove acceptance testing to regain schedule time.
\end{itemize}

Justify the strongest action using at least three indicators:

\vspace{10mm}\blankline

\section{9. Self-Check and Feedback}

Consult this section only after attempting all checkpoints.

\begin{checkbox}
\textbf{Expected calculations:}
\begin{itemize}
  \item PV = USD 40,000 $\times$ 0.50 = \textbf{USD 20,000}
  \item EV = USD 40,000 $\times$ 0.38 = \textbf{USD 15,200}
  \item SV = USD 15,200 $-$ USD 20,000 = \textbf{$-$USD 4,800}
  \item CV = USD 15,200 $-$ USD 24,000 = \textbf{$-$USD 8,800}
  \item SPI = 15,200 $\div$ 20,000 = \textbf{0.76}
  \item CPI = 15,200 $\div$ 24,000 = \textbf{0.63}
\end{itemize}
\end{checkbox}

\textbf{Expected diagnosis:} The project is behind schedule and over budget. It has delivered only USD 0.76 of value for every USD 1.00 planned by this point and USD 0.63 of value for every USD 1.00 spent. Module completion, critical bugs, and integration exposure also show a serious quality and delivery risk.

\textbf{Strongest action:} Freeze non-critical scope and concentrate resources on integration, essential pilot modules, critical defects, and acceptance testing. Report the revised forecast and required decision transparently.

If your answer differed, identify whether the error came from the formula, arithmetic, sign, rounding, or interpretation:

\vspace{8mm}\blankline

\section{10. Assessment and AI Boundary}

Assessment: calculations (12), interpretation (8), integrated diagnosis (5), and corrective-action justification (5), totaling 30 points.

Complete calculations independently before using AI. If AI is permitted, use it to explain an error after comparing your work with the self-check. Never accept an unexplained result.

\vfill
\begin{center}\small\color{pmoTeal}Course portal: \href{https://maoude.github.io/teaching-ai-pmo/}{maoude.github.io/teaching-ai-pmo/}\end{center}
\end{document}
